\chapter{Resultados e Discussão}
\label{chap:resultados}

Os resultados desta seção foram separados por nível de evidência. Os números da
rodada canônica estão em \path{evidencias/final_2026/outputs} e podem ser
rastreados ao manifesto por hashes. Onde a execução universal não foi concluída,
o status é apresentado como inconclusivo. Essa distinção é especialmente
importante em verificação baseada em BMC, cuja conclusão depende do programa,
do domínio e do solver efetivamente executados \citeonline{gopinath2021verifying}.

\section{Visão geral dos experimentos}

\begin{table}[H]
\centering
\caption{Resumo da rodada canônica e da rodada estrutural corrigida.}
\label{tab:resumo-experimentos}
\begin{tabularx}{\textwidth}{p{0.25\textwidth}p{0.24\textwidth}p{0.18\textwidth}Y}
\toprule
Experimento & Pergunta & Status & Evidência \\
\midrule
Quantização & Qual a diferença Float32--Q8.8 em 10.000 estados? & Confirmado & \path{quantization_canonical.json} e CSV de extremos \\
Aproximação tanh & O erro da função usada é inferior a 3\%? & Não confirmado & Varredura encontrou máximo de 4,5749\% \\
Reprodução A-direita & O estado publicado realmente gera \(z<0\)? & Confirmado para estado fixo & Harness C e log ESBMC \\
Reprodução B & O estado publicado sai da faixa em um passo? & Confirmado para estado fixo & Harness C e log ESBMC \\
Limite de força C & A força fica em \([-10,10]\,\mathrm{N}\) para todo o domínio? & Universal confirmado com Z3 & Harness simbólico, Z3 v4.8.9 \\
A-direita universal & A direção correta vale em todo o domínio? & Inconclusivo & Z3 atingiu timeout de 12 s \\
Atividade dos neurônios & Há contraexemplo de ativação para cada neurônio? & L1: 24/24; L2: 17/24; 7 TIMEOUT & \path{corrected_dead_neurons.json} \\
Saturação das camadas & Há contraexemplo de pré-ativação negativa? & L1: 24/24; L2: 22/24; 2 TIMEOUT & \path{corrected_saturation.json} \\
Saída do controlador & Os dois sinais da saída foram refutados universalmente? & Não: FAILED e TIMEOUT & \path{corrected_saturation.json} \\
\bottomrule
\end{tabularx}
\end{table}

\begin{figure}[H]
\centering
\begin{tikzpicture}[node distance=9mm and 9mm]
  \node[pibicbox, text width=3.5cm] (question) {Pergunta central\\o controlador quantizado\\é verificável e confiável?};
  \node[pibicresult, below left=of question, text width=3.7cm] (numeric) {Fidelidade numérica\\média: 0,0778 N\\máximo: 7,7589 N\\inversão de sinal no pior estado};
  \node[pibicresult, below=of question, text width=3.7cm] (formal) {Propriedades formais\\força limitada: provada\\A-direita universal: timeout\\replays A e B: reproduzidos};
  \node[pibicresult, below right=of question, text width=3.7cm] (structure) {Estrutura da rede\\L1: 24/24 contraexemplos\\L2: timeouts parciais\\saída conjunta: inconclusiva};
  \node[pibicprocess, below=of formal, text width=5.4cm] (conclusion) {Conclusão equilibrada\\há evidência reprodutível de limites e falhas concretas,\\mas não há prova de segurança global do controlador};
  \draw[pibicarrow] (question) -- (numeric);
  \draw[pibicarrow] (question) -- (formal);
  \draw[pibicarrow] (question) -- (structure);
  \draw[pibicarrow] (numeric) -- (conclusion);
  \draw[pibicarrow] (formal) -- (conclusion);
  \draw[pibicarrow] (structure) -- (conclusion);
\end{tikzpicture}
\caption{Síntese visual dos resultados e da conclusão sustentada pelos dados.}
\label{fig:sintese-resultados}
\end{figure}

A Figura~\ref{fig:sintese-resultados} mostra a conexão entre as análises. A
quantização apresenta boa média, mas uma cauda de erro capaz de inverter a
ação; a verificação formal prova o limite de força, porém não resolve a
propriedade direcional universal; e a auditoria estrutural encontra
contraexemplos, mas também timeouts. Em conjunto, os resultados justificam uma
conclusão de verificabilidade parcial e reprodutível, não uma certificação de
segurança global.

\section{Fidelidade Float32 versus Q8.8}

A análise foi reexecutada com \texttt{numpy.random.RandomState(42)}, quatro
variáveis amostradas uniformemente nos intervalos declarados e 10.000 estados.
Os pesos Float32 e Q8.8 foram lidos dos arquivos exportados; seus hashes estão
em \path{outputs/quantization_canonical.json}. Os valores coincidem com o
relatório histórico dentro do erro de representação decimal.

\begin{table}[H]
\centering
\caption{Erro absoluto entre a saída Float32 e a saída Q8.8.}
\label{tab:fidelidade}
\begin{tabular}{lr}
\toprule
Métrica & Valor \\
\midrule
Número de estados & 10.000 \\
Escala Q8.8 & 256 \\
Erro médio & 0,0778007759 N \\
Mediana & 0,0390625 N \\
Percentil 95 & 0,2262005212 N \\
Percentil 99 & 1,0927201962 N \\
Erro máximo & 7,7589197857 N \\
Máximo relativo à faixa de 10 N & 77,5892\% \\
\bottomrule
\end{tabular}
\end{table}

O maior erro ocorreu no estado
\[
(-0{,}8222019265,-2{,}4912202090,0{,}0051761657,0{,}7778176171),
\]
para o qual a saída Float32 foi \(+6{,}4307947857\,\mathrm{N}\) e a saída Q8.8
foi \(-1{,}328125\,\mathrm{N}\). Portanto, há inversão do sinal da ação nesse
caso. O resultado não deve ser descrito como equivalência ou “fidelidade total”
ao Float32. A afirmação sustentada é mais restrita: o harness C e o runtime web
usam a mesma representação e as mesmas operações Q8.8.

Uma varredura de \(z=-10000\) a \(z=10000\), documentada no arquivo JSON da
aproximação, encontrou erro absoluto máximo de
0,0457493, isto é, 4,5749\% do intervalo unitário da tanh, em \(z=-281\).
Esse valor contradiz a afirmação anterior de erro inferior a 3\%. O erro relativo
próximo de zero é mal condicionado e não é utilizado como métrica principal.

\section{Propriedades formais do controlador}

\subsection{Limite de força}

O harness universal da propriedade C usa estados simbólicos no domínio da
Tabela~\ref{tab:dominio-q88}, a rede inteira Q8.8 e a aproximação de tanh da
implementação. Com Z3 v4.8.9, a asserção
\(-2560\leq F_Q\leq2560\) terminou em \texttt{SUCCESSFUL} após 3,5 s. O
resultado é uma prova para esse harness e para esse domínio, não uma prova sobre
o Float32 nem sobre estados fora dos limites declarados. A propriedade também é
esperada pela construção da saída limitada a \(\pm255\) unidades Q8.8, mas a
execução simbólica fornece o registro operacional correspondente.

O mesmo teste foi tentado com \texttt{--boolector}. O executável Windows
disponível respondeu que o solver não foi compilado; esse caso foi registrado
como \texttt{UNKNOWN}. Assim, não se deve atribuir a prova desta rodada ao
Boolector.

\subsection{Segurança em um passo e reprodução do contraexemplo}

O estado histórico publicado para a propriedade B é
\[
s_0=(-0{,}7539,-3{,}9219,-0{,}1836,-1{,}5234).
\]
Após a quantização, ele corresponde a
\((-193,-1004,-47,-390)\). No harness Q8.8, a força calculada é
aproximadamente \(-9{,}9609\,\mathrm{N}\), e a atualização linearizada resulta em
\(\theta_1=-54/256\), aproximadamente \(-12{,}09^\circ\), ultrapassando o
limite inteiro \(-53\). O harness independente
\path{property_b_safety_replay.c} teve veredicto \texttt{SUCCESSFUL} para a
asserção fixa \(\theta_1<-53\), em 0,313 s.

Essa execução confirma a reprodução da violação alegada para o estado fixado.
Ela não mostra que todos os estados são inseguros e também não garante que o
controlador falhará na dinâmica não linear completa. O que foi verificado é o
modelo angular linearizado, em um único passo e na aritmética Q8.8.

\subsection{Correção direcional e resultados inconclusivos}

O estado publicado para A-direita é
\[
s_0=(-0{,}0117,-4{,}9219,0{,}1094,0),
\]
ou \((-3,-1260,28,0)\) em Q8.8. A reprodução independente verificou
\(z<0\) para esse estado, contrariando a propriedade que exigia \(z>0\); o
harness fixo terminou \texttt{SUCCESSFUL} em 0,609 s.

Já o teste universal A-direita, sem bounds artificiais de pré-ativação, atingiu
\texttt{TIMEOUT} de 12 s com Z3. Isso não é prova nem refutação universal. Há
uma divergência documental: o texto de apresentação antigo classifica as duas
direções como timeout, enquanto o JSON histórico classifica A-direita e
A-esquerda como \texttt{FAILED}. A-direita possui estado completo e foi
reproduzida; A-esquerda contém apenas \texttt{z=27}, sem estado completo, e não
é considerada evidência reproduzível.

\section{Análises estruturais corrigidas}

As análises históricas foram refeitas com o grafo completo da rede. As cópias
originais estão no diretório \texttt{legacy}; os resultados e logs da rodada
corrigida estão nos diretórios \texttt{outputs} e \texttt{logs}. Em particular,
a camada 2 é calculada a
partir do mesmo vetor simbólico da camada 1, sem injetar ativações independentes,
e não há bounds de pré-ativação artificialmente menores que os bounds derivados
do grafo.

\begin{table}[H]
\centering
\caption{Resultados da rodada estrutural corrigida com Z3.}
\label{tab:estrutura-corrigida}
\small
\begin{tabularx}{\textwidth}{>{\raggedright\arraybackslash}p{0.24\textwidth}rrr>{\raggedright\arraybackslash}X}
\toprule
Propriedade & \(\texttt{FAILED}\) & \(\texttt{SUCCESSFUL}\) & \(\texttt{TIMEOUT}\) & Interpretação \\
\midrule
Ativação -- camada 1 & 24 & 0 & 0 & Cada neurônio tem um contraexemplo ativo \\
Ativação -- camada 2 & 17 & 0 & 7 & 7 casos permanecem inconclusivos \\
Não saturação -- camada 1 & 24 & 0 & 0 & Cada neurônio tem um contraexemplo com \(pre<0\) \\
Não saturação -- camada 2 & 22 & 0 & 2 & 2 casos permanecem inconclusivos \\
Saída \(z\geq0\) & 1 & 0 & 0 & Há contraexemplo para a propriedade universal \\
Saída \(z\leq0\) & 0 & 0 & 1 & A execução terminou por tempo \\
\bottomrule
\end{tabularx}
\end{table}

Na verificação de atividade, \texttt{FAILED} significa que o solver encontrou
um estado do domínio em que o neurônio está ativo; portanto, os dados sustentam
24/24 contraexemplos na camada 1 e 17/24 na camada 2. Não sustentam, por si
sós, que o neurônio esteja ativo em todos os estados. Na camada 2, os sete
\texttt{TIMEOUT} não podem ser classificados como ativos ou mortos.

Na verificação de saturação, a asserção foi \(pre\geq0\), isto é, saturação
positiva em toda a região. Os 24 contraexemplos da camada 1 e os 22 da camada 2
mostram que cada um desses casos possui pelo menos uma entrada com
\(pre<0\). Os dois timeouts da camada 2 não autorizam conclusão. Assim, não há
neurônio cuja saturação positiva tenha sido provada universalmente nesta rodada.

Para a saída, \(z\geq0\) foi \texttt{FAILED} com o estado
\((-2,-305,-27,-1184)\) em Q8.8, enquanto \(z\leq0\) atingiu
\texttt{TIMEOUT} de 4 s. A combinação não permite declarar a saída
responsiva em ambos os sinais nem provar uma saturação global; o resultado
correto é inconclusivo para a propriedade conjunta.

\section{Escalabilidade e reprodutibilidade}

Os tempos confiáveis desta rodada são os dos harnesses gerados e dos logs
correspondentes: 0,609 s e 0,313 s para as reproduções fixas, 3,5 s para a
propriedade C universal com Z3 e 12 s até timeout para A-direita universal.
Na rodada estrutural corrigida, cada verificação foi limitada a 8 s na análise
de atividade e a 4 s na análise de saturação; esses limites estão registrados
nos JSONs. Os tempos históricos de 1888, 2722, 90 e 24 não são reutilizados,
pois o script atual não os mede nem os grava.

O fluxo demonstra uma diferença relevante entre detectar uma violação concreta
e provar uma propriedade para um domínio inteiro. A primeira pode ser reduzida a
um harness pequeno e reproduzível; a segunda exige solver disponível, modelo
sem restrições indevidas e tempo suficiente para explorar as combinações das
ReLUs. Essa distinção é parte do resultado tecnológico do projeto.

\section{Discussão integrada, limitações e ameaças à validade}

A rodada confirma três contribuições úteis: (i) uma comparação quantitativa e
rastreável entre Float32 e Q8.8; (ii) a reprodução independente de dois estados
que demonstram comportamentos alegados; e (iii) uma prova simbólica do limite de
força com Z3. Ao mesmo tempo, ela identifica limites que precisam aparecer no
relatório: a quantização pode inverter a ação em casos extremos, a tanh
aproximada tem erro maior que o anunciado, a propriedade de segurança tem
horizonte de um passo e o teste direcional não escala na configuração curta.
A rodada estrutural corrigida acrescenta contraexemplos de ativação e não
saturação para a maior parte dos neurônios, mas mantém sete casos de atividade
e dois de saturação como timeouts; portanto, não há base para uma afirmação
universal sobre os 48 neurônios.

Há ainda ameaças de modelagem: a dinâmica de B é linearizada, o domínio angular
inteiro é ligeiramente menor que a imagem completa do arredondamento de 12 graus,
e não foi realizada nesta rodada uma reprodução da trajetória não linear completa.
Por fim, o histórico de treinamento contém 1.670 episódios e média final de 100
episódios igual a 472,19, enquanto o gerador web histórico grava 500 e 475,0;
esses metadados devem ser reconciliados antes de serem usados como resultado.
